% r03_fill_K3b: proof of (K_0), part V: the comparison of the effective actions in total form, at h = 0 and with a wall.

\subsubsection*{The comparison of the effective actions at equal spacing}

\begin{theorem}[equal-spacing comparison, total form]\label{rot:thm:A}
Let $T\in\{T_1,T_2\}$, $G$ a real smooth $P$-periodic $2$-form, $z\in\R$, and $r=r(k)$ as in \Cref{rot:lem:r}. Then for $k\ge k_0(zG)$,
\[
\sup_{\supp\chi_r}\big|\hs^{(r)}_T(A_r,\phi_r;zG,0)-S^{(r)}_T(A_r,\phi_r;zG,0)\big|\le\vartheta_k|T|,\qquad
\vartheta_k\defeq C\big(5^{-\gamma r}\delta_r^{-\beta}(1+k)^p+\delta_r^\sigma\big)\to0,
\]
every $c$-number included: there is no vacuum constant and no source-linear constant at $h=0$.
\end{theorem}

\begin{theorem}[the vacuum constant]\label{rot:thm:V}
Let $v_k(T)\defeq S^{(r)}_T(0,0;0)-\hs^{(r)}_T(0,0;0)$. Then: \textup{(a)} $v_k=v^{\rm G}_k+v^{\rm I}_k$ with $v^{\rm G}=\ln N^Y_r-\ln\hat N_r+\ln N^Y_r(0)-\ln\hat N_r(0)$ and $|v^{\rm G}_k(T)|\le12.5\,a\,5^{3k-5r}|T|$; \textup{(b)} $E_1$ cancels exactly in $v^{\rm I}$, and $v^{\rm I}_k(T)=|T|\Delta p_k+\varpi_k(T)$ with $\Delta p_k$ a difference of infinite-lattice vacuum densities, the same for both tori, and $|\varpi_k(T)|\le C\eps^{-q}e^{-c\delta_r^{-1}}$; \textup{(c)} the order-$\lambda$ part obeys $|v^{\rm I}_\lambda|\le C\lambda\eta^{1/2}\delta_r^{-2}|T|+|\varpi_k|$; \textup{(d)} the order-$\lambda^2e^0$ part obeys $|v^{\rm I}_{\lambda^2}|\le C\lambda^2\delta_r^{-2}(\eps/\delta_r)^{1/4}\log(1/\eps)^2|T|+|\varpi_k|$; \textup{(e)} $|v^{\rm I}_k(T)|\le\vartheta_k|T|$. Hence $v_k(T)\to0$ on each torus.
\end{theorem}
\begin{proof}
(a) \cite[(3.89)]{K1} gives $E_0=\ln Z_r+\ln Z_r(0)+\ln N_r+\ln N_r(0)$; inserted in \cite[(3.14)]{K1} at $A_r=\phi_r=0$ this leaves $\ln N_r+\ln N_r(0)+P^{(r)}(0)$ for each flow, with $\ln N_r=-\frac32\sum_{p\in T^{(r)*}}\ln(\Delta^{(r)}_{\lmu}(p)/2\pi)$ and $\ln N_r(0)=-\frac N2\sum_p\ln(\Delta^{(r)}_m(p)/2\pi)$ over the dual of the common coarse torus. By \Cref{rot:prop:fc}, $|\ln(A^X(p)/A^Y(p))|\le0.20\,a\eta^2$ at every $p$ and every mass; sum over $|T|\delta_r^{-3}$ momenta and $3+N=5$ components with the factor $\frac12$; $\eta^2\delta_r^{-3}=25\cdot5^{-2r}5^{3(k-r)}$. (b) $E_1=|T|e_1(\eps)$ is the same $c$-number for both lattices and passes through every block-spin step unchanged. The remaining vertices at zero coarse fields are \cite[(3.18), (3.19), (3.22), (3.23)]{K1} and the products (3.25)$\times$(3.26), (3.26)$\times$(3.26), with fluctuation kernels as lines; each torus kernel is the $P$-periodisation of its infinite-lattice kernel, decaying at rate $\delta_*/\delta_r$ in unit scale, with prefactor $\le C\eta'^{-3}$; expanding every torus line as infinite-lattice line plus winding remainder, a graph with $n\le2\bar n$ lines differs from its infinite-lattice value by $\le C^n\eta'^{-3n}|T|\delta_r^{-3}\sum_{\gamma\ne0}e^{-\delta_*|\gamma|/2}\le C\eps^{-q}e^{-c\delta_r^{-1}}$, the shortest period being $50MJ\ge50\delta_r^{-1}$ in unit scale. The infinite-lattice values are $|T|p^F_k$ by $Y_{k-r}$-periodicity of both flows. (c) The order-$\lambda e^0$ vertices are (3.18) and the order-$\lambda$ part of (3.19), and $\dm_1=-4(N+2)\lambda C(0)$ (the vanishing of the one-particle-irreducible two-point function at zero momentum at order $\lambda$, \cite[(3.33)]{K1}). With Gaussian moments, $P^F(0)|_\lambda=\lambda N(N+2)\sum_x[(C(0)-G^F_T(x,x))^2-C(0)^2]$; $C(0)-G^F_r(x,x)=\Gamma^F(x,x)$ by \Cref{rot:lem:Gam}(i); the terms $-C(0)^2$ give equal totals; re-index the $Y$-sum by $\tau$ and use \Cref{rot:lem:Gam}(ii),(iii): $|\Gamma^X(x',x')^2-\Gamma^Y(\tau x',\tau x')^2|\le2C_1C_\Gamma\eta^{1/2}\delta_r^{-2}$. (d) With $V_1\defeq\lambda\sum{:}|\phi|^4{:}_C$ there are no tadpoles at order $\lambda^2$; $\frac12\dm_2=16(N+2)\lambda^2c_3$, $c_3=\sum_yC(y)^3$; writing ${:}F{:}_C=e^{-\frac12\Delta_\Gamma}{:}F{:}_G$ with $\Gamma=C-G$, the order-$\lambda^2e^0$ part of $v^{\rm I}$ is the $Y$-minus-$X$ difference of $N(N+2)\lambda^2[16T_1-24T_2+16T_3-4T_4-4(N+2)T_5]$ with $T_1=\sum_{x,y}C_{xy}^3(\Gamma_{xy}-\Gamma_{xx})$, $T_2=\sum C^2_{xy}\Gamma^2_{xy}$, $T_3=\sum C_{xy}\Gamma^3_{xy}$, $T_4=\sum\Gamma^4_{xy}$, $T_5=\sum\Gamma_{xx}\Gamma_{yy}(C_{xy}-\Gamma_{xy})^2$, up to $O(\eps^{-q}e^{-c\delta_r^{-1}})$ (\Cref{rot:prop:V2}); the pure-$C$ terms have equal totals. Each difference is bounded, in the three regions $|u|\le\rho\defeq(\eps\delta_r)^{1/2}$, $\rho<|u|\le R_\eps\defeq(40/m)\ln(1/\eps)$ and $|u|>R_\eps$, by the pointwise bounds $|C_{xy}|\le C_0(|u|+\eps)^{-1}$, $|\Gamma_{xy}|\le C'_1\min(\delta_r^{-1},(|u|+\eps)^{-1})$, $|\Gamma_{xy}-\Gamma_{xx}|\le C_2\delta_r^{-1}(|u|/\delta_r)^{1/2}$, the matched-pair bounds of \Cref{rot:cor:C} and \Cref{rot:lem:Gam}(iii), and Combes--Thomas decay beyond $R_\eps$; the result is $\le C\delta_r^{-2}(\eps/\delta_r)^{1/4}\ln(1/\eps)^2$ per unit volume. (e) Group the vacuum graphs, with $E_1$ expanded in each flow's slices (\Cref{rot:lem:ctsame}), into renormalised vacuum families in per-flow normal form; all lines are $A=0$ lines on the infinite lattice; the vacuum graphs reduce to families with at least two legs, to which \Cref{rot:prop:cp} applies per coarse site; multiply by $|T|\delta_r^{-3}$ and use \Cref{rot:lem:r}.
\end{proof}

\begin{proposition}[the order-$\lambda^2$ identity]\label{rot:prop:V2}
The decomposition of $v^{\rm I}_{\lambda^2}$ into $T_1,\dots,T_5$ in the proof of \Cref{rot:thm:V}(d) is an exact identity; in particular the counterterm $\dm_2$ enters only through $T_1$, in the form $\sum C^3[\Gamma(x,y)-\Gamma(x,x)]$.
\end{proposition}
\begin{proof}
Wick's theorem for the two Gaussians with covariances $C$ and $G=C-\Gamma$ at second order in $V_1$, with ${:}|\phi|^4{:}_C$ re-expanded in $G$-normal order: the connected contractions of two quartic vertices with $n$ lines of type $C$ and $4-n$ of type $\Gamma$, $n=0,\dots,4$, give $T_4,T_3,T_2,T_1$ and the pure-$C$ term, with the combinatorial weights $4$, $16$, $24$, $16$; the counterterm $\frac12\dm_2|\phi|^2$ contracted with the quartic vertex gives the $-c_3\sum\Gamma(x,x)$ half of $T_1$; the two one-loop pieces of the re-expansion contracted together give $T_5$. The form is the one in which each telescoping term is small (\Cref{rot:rem:tele}): with an unexpanded counterterm the intermediate terms would be of order $\lambda^2|T|\delta_r^{-1}$ and small only because $\sum_y\eps^3(C^{X3}-(C^Y\circ\tau)^3)=0$.
\end{proof}

\begin{remark}[why manifest forms]\label{rot:rem:tele}
The comparison must be applied to families in which every counterterm is expanded in the flow's own lines, and the same applies to Ward-identity cancellations, which hold for pure slice assignments only: King's ``replace one by one every factor'' \cite[p.~665]{K1} is correct on such manifest forms only (\Cref{rot:lem:ctsame}).
\end{remark}

\begin{proof}[Proof of \Cref{rot:thm:A}]
Fix $T$ and work on $\supp\chi_r$, with the bounds \cite[(3.2)]{K1}: $|A_r|\le p_r(\lmu\delta_r)^{-1}$, $|\partial A_r|\le p_r$, $|\phi_r|\le p_r(\lambda\delta_r)^{-1/4}$, $|D_{A_r}\phi_r|\le p_r$, $p_r\defeq p(\delta_r)$. Write both effective actions in the form of \Cref{rot:def:eff}, with the curvature source carried by the Cameron--Martin background (\Cref{rot:rem:CMrg}), and compare term by term; every bound is per coarse site, times $|T|\delta_r^{-3}$ sites.
\begin{enumerate}[leftmargin=*]
\item \emph{Photon quadratic form.} Three copies of the scalar coarse operator at mass $\lmu$, and $|\frac12\ip{A_r}{(\hat\Delta^{(r)}-\Delta^{(r)})A_r}|\le0.3a^2\eta^2p_r^2(\lmu\delta_r)^{-2}|T|\delta_r^{-3}$ by \Cref{rot:prop:fc} (King's (3.92) at equal spacing).
\item \emph{Normalisations.} \Cref{rot:thm:V}(a).
\item \emph{Scalar quadratic form at $A^{(r)}=0$.} As (1). Its $A^{(r)}$-dependent part has the form of a graph in $P^{(r)}$ with two external scalar fields \cite[p.~669]{K1} and is treated with (5).
\item \emph{The background determinant} $B3^F\defeq-\ln[Z^F_r(A^F)Z^F_r(0)^{-1}]$. By the scale factorisation of the determinant, valid for both flows by \Cref{rot:lem:tile}(iii),(iv), $B3^F=\frac12\sum_{j<r}\sum_x[g^F_{j,x}(A^F)-g^F_{j,x}(0)]$ with $g^F_{j,x}=\Tr_N[\ln K^F_j](x,x)$; the depths $j\le1$ are bounded per flow (per-site gauge and charge-conjugation parity), at cost $\le CNe^2p_r^2\delta_r^{-2}5^{-r}|T|$; the depths $j\ge2$ live on the common lattice $Y_{k-j}$ and are compared per site by an exact matched re-indexing, the telescoping of the resolvents along a cascade of exponential weights, a kernel sweep with the kernel bounds of \Cref{rot:prop:BD}(a) and \cite[Thm.~3.3]{K1}, and gradients compared as vectors, giving, for every $s_0>0$, $|B3^X-B3^Y|\le C_{s_0}Ne^2p_r^c(1+k)^p\delta_r^{-2}(5^{-\theta_Dr}+\delta_r^{s_0})|T|+C\eps^{-q}e^{-c/\delta_r}\le\vartheta_k|T|$, $\theta_D=\alpha/4$, by \Cref{rot:lem:r}. King's \cite[(3.94)--(3.98)]{K1} is not used: the per-$p$ middle line of (3.98) is false as printed, its seagull and bubble classes each diverging like $h^{-1}$ and cancelling only between classes.
\item \emph{The interaction $P^{(r)}$.} For each flow apply King's reduction \cite[(3.43)--(3.57)]{K1} in the form of \Cref{rot:lem:CB}: the constant-background gauge transformation (\Cref{rot:rem:frameav}); the localisation in universal-cover form; the expansions (3.47)--(3.54) around $\vartheta A^{(r)}=0$; the contour phases (\Cref{rot:conv:C}(ii), \Cref{rot:lem:shift}). After the reduction the background sits in the explicit factors of (3.55) and in the special lines $\delta^{\rm im}$, $\delta^{\rm bg}$, and every compared line is an $A=0$ line of \Cref{rot:cor:unif}; the contour vertices (3.24)--(3.26) and the phases (3.46) are compared by \Cref{rot:prop:BD}(b), with relative gains $\eta^{1/2}$ and $\eta'$ and the prefactor $p_r^2\delta_r^{-1/2}$; the explicit factors $A_r(w_1)\cdots A_r(w_q)$ are the same coarse field for both flows; the special lines, the terms with a vertex outside the vertex region and the Taylor remainders are bounded per flow by \Cref{rot:lem:CB}(2). Group the graphs into renormalised families in per-flow normal form; this is possible for both flows with the same combinatorics because each flow satisfies the identities used: counterterms (\Cref{rot:lem:ctsame}), inversion about coarse sites (\Cref{rot:conv:C}(i)), gauge covariance, and \cite[(3.87)]{K1}, an identity for each lattice's own full propagator. \Cref{rot:prop:cp} bounds each family difference by $Cp_r^2\delta_r^{-1/2}\eta'^{\theta}(1+k)^p$ times \cite[(1.33)]{Ba3}; the localisation sums give $C$ per coarse site, the field norms on $\supp\chi_r$ give at most $Cp_r^{\bar n}\delta_r^{-\beta_1}$ with $\beta_1$ fixed by \cite[(3.2)]{K1} and $\bar n$, and the number of families depends on $\bar n$ only.
\end{enumerate}
Summing, $|\hs^{(r)}-S^{(r)}|\le C(\eta'^{\theta}\delta_r^{-\beta}(1+k)^p+\delta_r^\sigma)|T|$ with $\eta^2\le25\eta'^\theta$ and the prefactor $p_r^2\delta_r^{-1/2}$ absorbed in $\delta_r^{-\beta}$ ($\beta\ge\beta_1+\frac12$ in \Cref{rot:lem:r}); with \Cref{rot:lem:r} this is $\vartheta_k|T|\to0$. The vacuum $c$-number is \Cref{rot:thm:V}(e). On $T_1$ the $Y$-flow lives on a torus skew in its own frame; this enters only through periodisation at $A=0$ and deck descent at $A\ne0$ (\Cref{rot:lem:F}).
\end{proof}

\subsubsection*{Walls: the source vertex}

\begin{lemma}[the $h$-census]\label{rot:lem:S1}
In King's degree (\cite[(3.30)--(3.31)]{K1}; $d=3$: a scalar end counts $-\frac12$, a differentiated end $-\frac32$, each factor $\eta$ counts $+1$, $D(G)=\sum_vD_G(v)-3$) the $h$-vertex \cite[(3.21)]{K1} has $D=3-\frac12\cdot2=2$ with both legs internal. Inserting one $h$-vertex into a scalar line of a graph $G_0$ gives $D(G_0+h)=D(G_0)+2$. Field-dependent $h$-free families have $D\ge-1$ with one exception, the current tadpole (one vertex \cite[(3.22)]{K1} with $m=1$ on a scalar loop, the photon leg external), which has $D=-2$: its value at $A=0$ vanishes ($\Tr q=0$), and in a background every insertion adds an external photon leg and raises the degree by $1$. So every non-vanishing field-dependent $h$-family has $D\ge1$. The only $h$-families with $D\le0$ are field-independent and linear in $h$: the free tadpole ($D=-1$; its residue at zero coarse fields, $-N\sum\eta^3h\,\Gamma^F(x,x)$, is Term~0 of \Cref{rot:prop:CK}), the order-$\lambda$ figure-eight with its $\dm_1$ partner, and the order-$e^2$ exchange and seagull tadpoles with their $\dm_2$ partner ($D=0$). Call this set $\cT_h$. The tree-level term $\delta_r^2\sum\eta^3h|\phi^{(r)}|^2$ has $D=0$ but no internal line; it is compared by \Cref{rot:lem:H} directly and gives $O(|T|\eps\delta_r^{-2}p_r^2)\to0$. The $X$/$\tau Y$ difference of the $h$-vertex is a smooth vertex function sampled at $x'$ and $\tau x'$ (\Cref{rot:lem:vert}(ii)), an explicit factor $\eta'$; $\sum_i\phi_i^2$ is frame-free.
\end{lemma}
\begin{proof}
Degree arithmetic from the vertex list \cite[(3.18)--(3.26)]{K1} (the seagull of order two included) and the minimal degree $-1$ of the field-dependent $h$-free families read from Ba{\l}aban's list \cite[(2.4)]{Ba3}; the exhaustive enumeration of the $h$-linear graphs of order $\le2$ in the couplings, with the merged counterterm vertices, gives $\cT_h$. For a one-particle-reducible family a bridge line adds $-1+3=+2$ to the degree, so its degree is at least the minimum over its one-particle-irreducible parts.
\end{proof}

\begin{lemma}[the source-linear constant is the same on both tori]\label{rot:lem:C}
Let $f^F(x)$ be the infinite-lattice one-point density of $\cT_h$ along flow $F$ (one $h$-vertex at $x$, all lines infinite-lattice kernels); it is $Y_{k-r}$-periodic. Let $\bar f^F$ be its cell average and $c_k\defeq-(\bar f^X-\bar f^Y)$, a number independent of $T$. For $T\in\{T_1,T_2\}$ the $\cT_h$-part of $\hs_T-S_T$ at zero coarse fields equals $-c_k\sS_{\Xk}(h)+\rho_T(h)$ with $|\rho_T(h)|\le C(\eps^{-q}e^{-c\delta_r^{-1}}+C_N\eps^{-q}\delta_r^N)\|h\|_{C^N}\to0$. Moreover $|c_k|\le C\eps^{-q}$.
\end{lemma}
\begin{proof}
Periodisation as in \Cref{rot:thm:V}(b), with one $h$-vertex summed over $T$. Write $x=z+t$ with $z\in Y_{k-r}/P$ and $t$ in a cell; the Riemann sum of the smooth periodic $h$ over the coarse lattice is exact to $O(C_N\delta_r^N)$, so $\sum_x\eps^3h(x)f^F(x)=\bar f^F\int h+O(\|f^F\|_\infty\delta_r^N)$, with $\|f^F\|_\infty\le C\eps^{-q}$ since each line is $\le C\eta'^{-3}$; and $\sS_{\Xk}(h)=\int h+O(C_N\eps^N)$ (\Cref{rot:lem:lsum}). Take $N$ large: with $\delta_r=5^{-\lfloor\vr k\rfloor}$, $N>(3+q-3\vr)/\vr$ suffices.
\end{proof}

\begin{theorem}[comparison with a wall]\label{rot:thm:Asrc}
For $T\in\{T_1,T_2\}$ and $s=(zG,h)$ real, on $\supp\chi_r$: $\big|\hs^{(r)}_T(\cdot;s)-S^{(r)}_T(\cdot;s)+c_k\sS_{\Xk}(h)\big|\le\vartheta_k(s)|T|\to0$, with $c_k$ of \Cref{rot:lem:C}.
\end{theorem}
\begin{proof}
Split $\hs(G,h)-S(G,h)=[\hs(G,0)-S(G,0)]+[\Delta_h\hs-\Delta_hS]$. The first bracket is \Cref{rot:thm:A}. Since \cite[(3.17)]{K1} depends on $h$ only through the $h$-vertex, the second bracket consists of the $h$-families of the two Taylor polynomials plus Taylor remainders bounded per flow by $C\delta_r^\sigma|T|$ \cite[p.~665]{K1}. The $\cT_h$ part is $-c_k\sS_{\Xk}(h)+o(1)$ by \Cref{rot:lem:C}; its $G$-dependence through the Cameron--Martin background is an $h$-family with background insertions, which by gauge invariance of each flow acts through the curvature of the background and so has positive degree. Every other $h$-family has $D\ge1$ (\Cref{rot:lem:S1}); group them in per-flow normal form with the same combinatorics in both flows, telescope element by element (lines: \Cref{rot:cor:unif}; external fields: \Cref{rot:lem:H}; the $h$-vertex: \Cref{rot:lem:vert}(ii)), apply \Cref{rot:prop:cp}, and sum as in the proof of \Cref{rot:thm:A}. The exponents of \Cref{rot:lem:r} are source-independent (\cite[Thm.~3.4]{K1}: ``$C$ depends on $g,h$'').
\end{proof}
